% irreversible.tex — Irreversible substitution models and κ > 1 regime
% To be \input from lhopital-limits.tex

%%%%%%%%%%%%%%%%%%%%%%%%%%%%%%%%%%%%%%%%%%%%%%%%%%%%%%%%%%%%%%%%%%%%%%%%
\subsection{\texorpdfstring{Irreversible Models and the $\insrate > \delrate$ Regime}{Irreversible Models and the lambda > mu Regime}}
\label{sec:irreversible}

\paragraph{The $\kappa > 1$ regime.}
When $\insrate > \delrate$, the ratio $\kappa = \insrate/\delrate > 1$
and the geometric$(\kappa)$ prior on ancestor length is not
normalizable: $\sum_{i=0}^\infty \kappa^i(1-\kappa)$ diverges.
Therefore the joint likelihood $P(x,y)$ does not exist.
However, the ancestor length-conditioned likelihood
$P_{\mathrm{cond}}(x,y \mid |x|)$
is well-defined for all $\kappa > 0$.
The transition parameters $\alpha$, $\beta$, $\gamma$ are
continuous functions of $(\insrate,\delrate,\evoltime)$
with no singularity at $\insrate = \delrate$ or $\insrate > \delrate$.
The BDI sufficient statistics $\mathbb{E}[B]$, $\mathbb{E}[D]$,
$\mathbb{E}[S]$ remain well-defined, and the conditioned
M-step~\eqref{eq:mstep-cond} applies unchanged:
it is unconstrained in $\insrate$ and $\delrate$ and permits
$\hat\insrate > \hat\delrate$ if the data support it.

\paragraph{Irreversible substitution models.}
For completeness, our implementations support dropping the
reversibility assumption for the substitution CTMC.
The integrals described in Section~\ref{sec:bridge-expectations} and in~\cite{HolmesRubin2002}
exploit the symmetric factorization of $\revsub$
using $\eqm$-weighted similarity, which yields real eigenvalues.
In the irreversible case, $\revsub$ is a general rate matrix
with no detailed-balance constraint.
The eigendecomposition $\revsub = V\,\operatorname{diag}(d)\,V^{-1}$
may have complex eigenvalues.
Assuming $\revsub$ is diagonalizable, the spectral formulae for the integrals of Section~\ref{sec:bridge-expectations} carry over,
with $V^{-1}$ replacing $V^\top$ and complex arithmetic throughout:
\[
I^{ab}_{ij}(\evoltime) =
\sum_{k,l} V_{ak}\,V^{-1}_{ki}\,
\frac{e^{d_k\evoltime} - e^{d_l\evoltime}}{d_k - d_l}\,
V_{jl}\,V^{-1}_{lb}
\]
with the convention $(e^{d_k\evoltime} - e^{d_l\evoltime})/(d_k - d_l) = \evoltime\,e^{d_k\evoltime}$
when $d_k = d_l$.
The result is real (complex eigenvalues come in conjugate pairs).

\paragraph{Fixed-point iteration for $\eqm$.}
In the reversible case, the equilibrium distribution $\eqm$ is
constrained to be the stationary distribution of $\revsub$ by
by detailed balance, and may either be fixed a priori or updated jointly with the reversible rate parameters during the M-step.
In the irreversible case, $\eqm$ must be recomputed as the
left eigenvector of $\revsub$ (with eigenvalue~0) after each
M-step update of $\revsub$ from the bridge-expectation sufficient statistics.
This gives an ECM (Expectation-Conditional-Maximization)
scheme~\cite{MengRubin1993}:
the E-step uses the current $(\revsub, \eqm)$;
the M-step updates $\revsub$, then $\eqm$ is recomputed from
the updated $\revsub$.
Standard monotonicity results apply providing each conditional maximization step is well-defined.

\paragraph{Differentiability.}

Standard autograd libraries provide general eigendecomposition and matrix-exponential primitives (with Fréchet derivatives).
The eigendecomposition route is convenient for EM-style closed-form updates, while gradient-based optimization is often more naturally based on the matrix exponential. We do not rely on differentiating through eigenvectors here.
